% ============================================================
% The General Neutrality Theorem
% ============================================================

\documentclass[12pt]{amsart}

\usepackage{amsmath,amssymb,amsthm}
\usepackage{booktabs}
\usepackage{microtype}
\usepackage{url}

\newtheorem{theorem}{Theorem}
\newtheorem{lemma}[theorem]{Lemma}
\newtheorem{proposition}[theorem]{Proposition}
\newtheorem{corollary}[theorem]{Corollary}
\theoremstyle{definition}
\newtheorem{definition}[theorem]{Definition}
\theoremstyle{remark}
\newtheorem{remark}[theorem]{Remark}

\title{The General Neutrality Theorem}

\author{Alexander S.\ Petty}
\email{alexander.petty@gmail.com}
\date{September 2025 (revised April 2026)}
\subjclass[2020]{11A63, 11N05, 11M06}

\begin{document}
\begin{abstract}
The collision invariant $S_{\ell}$ satisfies the
reflection identity $S(a) + S(m{-}a) = -1$ on
$(\mathbb{Z}/m\mathbb{Z})^{\times}$. The companion
paper proved neutrality at $q = 3$: a specific
mod-$3$ class has mean $-1/2$. This paper proves the
general version: for every odd prime $q$ not dividing
$b$, the reflection fixes exactly one residue class
modulo $q$, and that class has mean $-1/2$ on the
combined CRT period modulo $mq$.

These constraints can be imposed simultaneously on the
combined CRT period for finitely many primes. By
Mertens' third theorem, the fraction of lifted units
that avoid neutrality at every prime $q \le Q$ tends
to zero as $Q \to \infty$.

The centered prime harmonic sum
$F^{\circ}(s) = \sum \hat{S}^{\circ}(\chi)\, P(s, \chi)$
decomposes over odd characters. Computation of the
products $\hat{S}^{\circ}(\chi) \cdot P(s, \chi)$
reveals that they have mixed signs: at $s = 1/2$ in
base~$10$, thirteen of twenty are positive and seven
are negative. The positive terms outweigh the negative
in magnitude, not in count. This pattern persists
across all prime bases tested up to $b = 37$, with the
positive fraction drifting toward $50\%$ as the number
of characters grows, while the net sum remains positive.

The data therefore points toward a collective
phenomenon. Convergence of $F^{\circ}$ would constrain
the joint zero distribution
of the odd $L$-functions modulo $m$, not each
$L$-function individually.
\end{abstract}
\maketitle

% ============================================================
\section{Introduction}

The collision invariant $S_{\ell}$ is a function on
$(\mathbb{Z}/m\mathbb{Z})^{\times}$ with
$m = b^{\ell+1}$, satisfying the reflection
identity~\cite{paperA}
\[
S(a) + S(m{-}a) = -1.
\]
The companion paper~\cite{paperB} proved convergence of
the centered prime harmonic sum $F^{\circ}(s)$ at
$s = 1$ and identified the neutrality theorem at
$q = 3$ as the mechanism controlling penetration
below~$s = 1$.

This paper extends the neutrality theorem to all odd
primes and examines the sign structure of the character
decomposition. The results separate the behavior of
$F^{\circ}$ from the simplest positivity route to the
Generalized Riemann Hypothesis: the mixed-sign
computations indicate that convergence of $F^{\circ}$
need not be a test of each $L$-function individually,
but rather a collective constraint on the family.

% ============================================================
\section{The General Neutrality Theorem}

\begin{definition}
Let $q$ be an odd prime with $q \nmid b$, and let
$m = b^{\ell+1}$. For a residue class $k \bmod q$,
define the lifted class
\[
\mathcal U_{m,q}(k)
=
\{1 \le a \le mq : \gcd(a,m)=1,\ a \equiv k \pmod q\}.
\]
Its lifted mean is
\[
\overline S^{(q)}_k
=
\frac{1}{\phi(m)}
\sum_{a \in \mathcal U_{m,q}(k)} S_\ell(a \bmod m).
\]
By the Chinese remainder theorem, each lifted class
contains exactly $\phi(m)$ elements.
\end{definition}

\begin{theorem}[General neutrality]\label{thm:general}
Let $q$ be an odd prime with $q \nmid b$, and let
$m = b^{\ell+1}$. Define
\[
k^{*} \equiv m \cdot \frac{q+1}{2} \pmod{q},
\]
the unique solution to $2k \equiv m \pmod{q}$. Then:
\begin{enumerate}
\item The reflection $a \mapsto m - a$ maps the class
$\mathcal U_{m,q}(k^{*})$ to itself modulo $mq$.
\item The lifted mean over this class equals $-1/2$:
\[
\overline S^{(q)}_{k^{*}} = -\frac12.
\]
\end{enumerate}
\end{theorem}

\begin{proof}
The reflection sends class $k$ to class
$m - k \pmod{q}$. The fixed class satisfies
$k \equiv m - k$, i.e., $2k \equiv m \pmod{q}$.
Since $q$ is odd, $2$ is invertible and the solution
$k^{*}$ is unique.

The reflection lifts to an involution on
$\mathcal U_{m,q}(k^{*})$ modulo $mq$: if
$a \in \mathcal U_{m,q}(k^{*})$, let $a^{\sharp}$ be
the unique representative in $[1,mq]$ satisfying
\[
a^{\sharp} \equiv m-a \pmod{mq}.
\]
Then $a^{\sharp} \equiv m-k^{*} \equiv k^{*} \pmod q$,
and $a^{\sharp} \bmod m \equiv -a \bmod m$, so
$a^{\sharp}$ is again coprime to $m$. Thus
$a^{\sharp} \in \mathcal U_{m,q}(k^{*})$.

Since $S_\ell$ is a function modulo $m$, the reflection
identity gives
\[
S_\ell(a \bmod m) + S_\ell(a^{\sharp} \bmod m) = -1.
\]
Summing over the involution pairs and dividing by
$|\mathcal U_{m,q}(k^{*})| = \phi(m)$ yields
$\overline S^{(q)}_{k^{*}} = -1/2$.
\end{proof}

\begin{corollary}[Paired classes]\label{cor:paired}
The remaining $q - 1$ classes split into $(q{-}1)/2$
pairs $\{k, \, m{-}k \bmod q\}$ swapped by reflection.
Their lifted means satisfy
\[
\overline S^{(q)}_{k} + \overline S^{(q)}_{m-k} = -1.
\]
\end{corollary}

\begin{corollary}[Simultaneous neutrality]\label{cor:simult}
For any finite set of odd primes $q \nmid b$, the
neutrality constraints from Theorem~\ref{thm:general}
hold simultaneously on the combined CRT period modulo
$m\prod q$.
\end{corollary}

% ============================================================
\section{The Neutral Fraction}

Fix a finite set $\mathcal Q$ of odd primes not
dividing $b$, and set
\[
M_{\mathcal Q} = \prod_{q \in \mathcal Q} q.
\]
Consider the lifted unit period
\[
\mathcal U_{m,\mathcal Q}
=
\{1 \le a \le m M_{\mathcal Q} : \gcd(a,m)=1\}.
\]
For each $q \in \mathcal Q$, let $k_q^*$ denote the
neutral class from Theorem~\ref{thm:general}. The
fraction of lifted units that avoid neutrality at
every prime in $\mathcal Q$ is exactly
\[
\prod_{q \in \mathcal Q} \frac{q - 1}{q}.
\]
Indeed, by the Chinese remainder theorem, the
conditions $a \equiv k_q^* \pmod q$ are independent
across $q \in \mathcal Q$. For each fixed $q$, exactly
one of the $q$ residue classes is neutral, so the
avoiding fraction is $(q-1)/q$, and the product is
exact on the combined CRT period. By Mertens' third
theorem, this product tends to zero as $\mathcal Q$
exhausts the odd primes not dividing $b$.

\begin{remark}
Each constraint fixes a class mean, not individual
values. The collision invariant is not determined by
neutrality. But it is far from generic: a random
function lifted to the CRT period would satisfy none
of these constraints.
\end{remark}

\begin{table}[h]
\centering
\begin{tabular}{rrrrrr}
\toprule
$b$ & $q$ & $k^{*}$ & $\sum S$ & count &
$2\sum{+}\text{count}$ \\
\midrule
$5$ & $3$ & $2$ & $-3$ & $6$ & $0$ \\
$5$ & $7$ & $2$ & $-2$ & $4$ & $0$ \\
$7$ & $3$ & $2$ & $-7$ & $14$ & $0$ \\
$7$ & $5$ & $2$ & $-4$ & $8$ & $0$ \\
$10$ & $3$ & $2$ & $-6$ & $12$ & $0$ \\
$10$ & $7$ & $1$ & $-3$ & $6$ & $0$ \\
$13$ & $3$ & $2$ & $-26$ & $52$ & $0$ \\
$13$ & $5$ & $2$ & $-16$ & $32$ & $0$ \\
$13$ & $7$ & $1$ & $-11$ & $22$ & $0$ \\
\bottomrule
\end{tabular}
\caption{The identity $2\sum S + \text{count} = 0$
(mean $= -1/2$) holds exactly for all tested bases
$b \le 20$ and primes $q \le 19$: $117/117$ cases.}
\label{tab:general}
\end{table}

% ============================================================
\section{The Sign Structure}

The antisymmetry theorem~\cite{paperB} gives
\[
F^{\circ}(s) = \sum_{\substack{\chi \bmod m \\
\chi(-1) = -1}} \hat{S}^{\circ}(\chi)\, P(s, \chi),
\]
a sum over non-trivial odd characters. Each product
$\hat{S}^{\circ}(\chi) \cdot P(s, \chi)$ is a complex
number whose real part contributes to $F^{\circ}(s)$.
If all these real parts were non-negative, convergence
of $F^{\circ}$ at $s = 1/2$ would require each
$P(s, \chi)$ to converge individually, equivalent to
the Generalized Riemann Hypothesis for the odd
$L$-functions modulo~$m$.

\begin{table}[h]
\centering
\begin{tabular}{rrrrrr}
\toprule
$s$ & positive & negative & $\sum_+ $ & $\sum_-$ &
net \\
\midrule
$1.0$ & $13$ & $7$ & $+0.47$ & $-0.39$ & $+0.08$ \\
$0.8$ & $15$ & $5$ & $+0.28$ & $-0.10$ & $+0.19$ \\
$0.6$ & $13$ & $7$ & $+0.75$ & $-0.31$ & $+0.44$ \\
$0.5$ & $13$ & $7$ & $+1.46$ & $-0.73$ & $+0.73$ \\
\bottomrule
\end{tabular}
\caption{Sign of
$\operatorname{Re}(\hat{S}^{\circ}(\chi) \cdot
P(s, \chi))$ across the $20$ odd characters
modulo~$100$. Base~$10$, lag~$1$, $348{,}488$ primes.
The products have mixed signs at every $s$.}
\label{tab:signs10}
\end{table}

The positivity condition fails. Seven of twenty products
are negative at $s = 0.5$ in base~$10$. The net sum is
positive because the positive terms have larger
magnitude, not because they are more numerous.

% ============================================================
\section{Universality of the Sign Structure}

\begin{table}[h]
\centering
\begin{tabular}{rrrrrr}
\toprule
$b$ & odd chars & positive & negative & $\%$pos &
net \\
\midrule
$3$ & $3$ & $3$ & $0$ & $100$ & $+0.62$ \\
$5$ & $10$ & $8$ & $2$ & $80$ & $+1.04$ \\
$7$ & $21$ & $13$ & $8$ & $62$ & $+1.81$ \\
$11$ & $55$ & $27$ & $28$ & $49$ & $+1.66$ \\
$13$ & $78$ & $50$ & $28$ & $64$ & $+3.45$ \\
$17$ & $136$ & $74$ & $62$ & $54$ & $+2.53$ \\
$19$ & $171$ & $99$ & $72$ & $58$ & $+5.11$ \\
$23$ & $253$ & $129$ & $124$ & $51$ & $+1.56$ \\
$29$ & $406$ & $224$ & $182$ & $55$ & $+5.32$ \\
$31$ & $465$ & $249$ & $216$ & $54$ & $+6.73$ \\
$37$ & $666$ & --- & --- & --- & $+8.25$ \\
\bottomrule
\end{tabular}
\caption{Sign split across prime bases at $s = 0.5$
($148{,}000$ primes per base). The positive fraction
drifts toward $50\%$ as the number of characters
grows. The net remains positive in every case.}
\label{tab:multi}
\end{table}

Two patterns emerge across all bases tested.

First, the positive fraction drifts toward $50\%$ as
$\phi(b^2)$ increases. At base~$3$ (three odd
characters) all products are positive. At base~$23$
($253$ odd characters) the split is $129$ to $124$.
The sign structure is not dominated by a majority.

Second, the net sum is positive at every base. The
positive terms win not by number but by magnitude:
the collision coefficients with the largest
$|\hat{S}^{\circ}(\chi)|$ tend to pair with
convergence-supporting (positive) products.

\begin{table}[h]
\centering
\begin{tabular}{rrcc}
\toprule
$b$ & odd chars & top-$1$ sign & top-$2$ sign \\
\midrule
$3$ & $3$ & $+$ & $+$ \\
$5$ & $10$ & $+$ & $+$ \\
$7$ & $21$ & $+$ & $+$ \\
$11$ & $55$ & $+$ & $+$ \\
$13$ & $78$ & $+$ & $+$ \\
$17$ & $136$ & $+$ & $+$ \\
$19$ & $171$ & $+$ & $+$ \\
$23$ & $253$ & $+$ & $+$ \\
$29$ & $406$ & $-$ & $-$ \\
$31$ & $465$ & $+$ & $+$ \\
$37$ & $666$ & $+$ & $+$ \\
\bottomrule
\end{tabular}
\caption{Sign of the product
$\hat{S}^{\circ}(\chi) \cdot P(s, \chi)$ at $s = 0.5$
for the characters with the two largest
$|\hat{S}^{\circ}|$. In $10$ of $11$ bases, the
dominant coefficient supports convergence. Base~$29$
is the sole exception, but the net remains positive.}
\label{tab:dominant}
\end{table}

At base~$29$, the two largest coefficients are
negative, but the third and fourth largest are positive
with stronger phase alignment, producing a net
contribution that overwhelms the negative leaders.
The structural integrity does not depend on any single
character.

% ============================================================
\section{The Anti-Correlation}

The sign structure raises a question: why do the
positive terms consistently outweigh the negative?
The answer is an anti-correlation between the
collision coefficients and the prime character sums.

\begin{table}[h]
\centering
\begin{tabular}{rrrr}
\toprule
$b$ & odd chars & $\rho$ (Pearson) & $\rho_s$ (Spearman) \\
\midrule
$3$ & $3$ & $+0.00$ & $-1.00$ \\
$5$ & $10$ & $-0.45$ & $+0.31$ \\
$7$ & $21$ & $-0.28$ & $+0.21$ \\
$11$ & $55$ & $-0.23$ & $-0.11$ \\
$13$ & $78$ & $-0.27$ & $-0.08$ \\
$17$ & $136$ & $-0.32$ & $-0.28$ \\
$19$ & $171$ & $-0.13$ & $-0.05$ \\
$23$ & $253$ & $-0.26$ & $-0.21$ \\
$29$ & $406$ & $-0.24$ & $-0.21$ \\
$31$ & $465$ & $-0.24$ & $-0.27$ \\
$37$ & $666$ & $-0.27$ & $-0.30$ \\
\bottomrule
\end{tabular}
\caption{Pearson and Spearman correlations between
$|\hat{S}^{\circ}(\chi)|$ and $|P(0.5, \chi)|$ across
odd characters. The Pearson correlation is negative at
every base: the collision invariant places its largest
Fourier weight on characters with moderate prime sums.
$148{,}000$ primes per base.}
\label{tab:anticorr}
\end{table}

The Pearson correlation between
$|\hat{S}^{\circ}(\chi)|$ and $|P(s, \chi)|$ is
negative at every prime base tested, ranging from
$-0.13$ to $-0.45$. Characters with large collision
coefficients tend to have moderate prime character
sums, and characters with large prime sums tend to
have small collision coefficients.

The prime character sum $|P(s, \chi)|$ measures a
character's sensitivity to the distribution of
primes in residue classes and, through the relation
$P = \log L - H$, to the zeros of the corresponding
$L$-function. Characters with large $|P|$ are those
most affected by nearby zeros.

The anti-correlation therefore suggests that the collision
invariant places its largest Fourier weight on the
characters \emph{least} sensitive to $L$-function
zeros. The digit geometry of $\lfloor br/p \rfloor$
appears to produce a spectral profile that is
misaligned with the zero-sensitive directions in
character space.

This points to a structural mechanism behind the
convergence of $F^{\circ}$. The sum may converge not
because the $L$-function zeros are absent, but
because the collision invariant is not looking at
them.

% ============================================================
\section{The Scramble Test}

The anti-correlation appears to be load-bearing: in
the scramble experiment, scrambling the assignment of
collision coefficients to characters destroys the
observed net.

\begin{proposition}\label{prop:scramble}
Let $\pi$ be a uniform random permutation of the odd
characters. The expected value of the scrambled sum
$\sum_{\chi} \hat{S}^{\circ}(\chi_{\pi}) \cdot
P(s, \chi)$ is
\[
\mathbb{E}\Bigl[\sum_{\chi}
\operatorname{Re}\bigl(\hat{S}^{\circ}(\chi_{\pi})
\cdot P(s, \chi)\bigr)\Bigr]
= \frac{1}{n}\operatorname{Re}\Bigl(
\sum_{\chi} \hat{S}^{\circ}(\chi) \cdot
\sum_{\chi} P(s, \chi)\Bigr),
\]
where $n$ is the number of odd characters. By
character orthogonality,
$\sum_{\chi \text{ odd}} \chi(p) =
(\phi(m)/2)(\delta_{p \equiv 1} -
\delta_{p \equiv -1})$, so
$\sum_{\chi} P(s, \chi)$ receives contributions only
from primes $p \equiv \pm 1 \pmod{m}$.
\end{proposition}

\begin{proof}
For a uniform random permutation $\pi$,
\[
\mathbb{E}[\hat{S}^{\circ}(\chi_{\pi(j)})]
= \frac{1}{n} \sum_{k} \hat{S}^{\circ}(\chi_k)
\]
for each $j$. Substituting and exchanging summation
gives the displayed formula. The character sum
identity is standard~\cite{iwaniec}.
\end{proof}

\begin{table}[h]
\centering
\begin{tabular}{rrrrrr}
\toprule
$b$ & odd & actual & $\mathbb{E}[\text{scr}]$ &
$\sigma[\text{scr}]$ & $z$ \\
\midrule
$7$ & $21$ & $+1.81$ & $-0.16$ & $0.92$ & $+2.1$ \\
$11$ & $55$ & $+1.66$ & $-0.06$ & $1.27$ & $+1.4$ \\
$13$ & $78$ & $+3.45$ & $-0.05$ & $1.43$ & $+2.5$ \\
$17$ & $136$ & $+2.53$ & $-0.05$ & $1.64$ & $+1.6$ \\
$19$ & $171$ & $+5.11$ & $-0.07$ & $1.66$ & $+3.1$ \\
$23$ & $253$ & $+1.56$ & $-0.00$ & $1.87$ & $+0.8$ \\
$29$ & $406$ & $+5.32$ & $+0.02$ & $2.09$ & $+2.5$ \\
$31$ & $465$ & $+6.73$ & $-0.00$ & $2.27$ & $+3.0$ \\
\bottomrule
\end{tabular}
\caption{The actual net $F^{\circ}(0.5)$ compared with
$10{,}000$ random permutations of the collision
coefficients. Proposition~\ref{prop:scramble} gives
the exact expected scrambled mean; in the tested
bases this mean is near zero. The actual net
is $1$--$3$ standard deviations above, indicating
the specific coefficient-to-character assignment is
relevant to the observed convergence pattern.}
\label{tab:scramble}
\end{table}

The $z$-scores are consistently positive and mostly
above $2$. A generic weighting of the same characters
would produce a net near zero. The collision
invariant's spectral profile is not merely different
from random: in the tested bases, it is specifically
oriented toward a positive net. The anti-correlation
between $|\hat{S}^{\circ}|$ and $|P|$ is the leading
candidate mechanism.

\begin{remark}
The scramble test is computational evidence, not a
proof. Proving that the collision invariant's Fourier
spectrum must be anti-correlated with the prime
character sum magnitudes is an open problem. It would
require understanding why the digit geometry of
$\lfloor br/p \rfloor$ produces a spectral profile
transverse to the zero-sensitive directions in
character space.
\end{remark}

% ============================================================
\section{Consequences}

\begin{remark}
The mixed-sign data shows that any implication from
convergence of $F^{\circ}$ at $s = 1/2$ to individual
zero-freeness of every $P(s,\chi)$ cannot proceed by a
positivity argument. If convergence of $F^{\circ}$ at
$s = 1/2$ holds, it is compatible with cancellation
between terms of opposite sign.
\end{remark}

The collision transform is therefore not a detector of
individual $L$-function zeros. If $F^{\circ}$
converges at $s = 1/2$, the statement is not that
every odd $L$-function modulo $m$ is zero-free for
$\operatorname{Re}(s) > 1/2$. The statement is that
the zeros, if any, are arranged so that their effects
cancel under the collision weighting. This is a
constraint on the geometry of the joint zero set,
not on its emptiness.

\begin{remark}
The ratio of positive to negative total magnitude is
approximately $2{:}1$ at $s = 0.5$ in base~$10$
(Table~\ref{tab:signs10}). Whether this ratio is
universal, base-dependent, or depth-dependent is an
open question.
\end{remark}

% ============================================================
\section{Remarks}

\subsection*{What is proved}

The general neutrality theorem is proved for all odd
primes $q$ not dividing $b$, using only the reflection
identity and the invertibility of $2$ modulo $q$.
Computational verification covers $117$ cases (all
bases $b \le 20$, primes $q \le 19$).

The sign structure, anti-correlation, and scramble
behavior are computational observations supported by
$11$ prime bases with up to $666$ odd characters each.
They are not theorems.

\subsection*{The structure at $q = 2$}

The prime $q = 2$ is excluded because $2$ is not
invertible modulo $2$. The parity structure of the
collision invariant depends on the parity of $b$
and requires separate analysis.

\subsection*{The arc}

The neutrality theorem at $q = 3$ controls the
cancellation mechanism that allows $F^{\circ}$ to
penetrate below $s = 1$~\cite{paperB}. The general
neutrality theorem shows this is the first instance
of a family of constraints at every prime scale. The
computed sign structure indicates that the convergence of
$F^{\circ}$ is not a GRH equivalence but a collective
phenomenon: the collision invariant creates a
structured weighting of $L$-functions in which
individual zeros can be absorbed by the cancellation
between terms of opposite sign.

% ============================================================
\begin{thebibliography}{99}

\bibitem{davenport}
H.~Davenport,
\emph{Multiplicative Number Theory},
3rd~ed., Springer, 2000.

\bibitem{iwaniec}
H.~Iwaniec and E.~Kowalski,
\emph{Analytic Number Theory},
AMS Colloquium Publications, vol.~53, 2004.

\bibitem{paperA}
A.~S.~Petty,
The collision invariant,
preprint, 2026.

\bibitem{paperB}
A.~S.~Petty,
The collision transform,
preprint, 2026.

% SUBMISSION GATE: replace this moving repository URL with a
% fixed release/tag/commit/checksum/DOI before submission.
\bibitem{nfield}
A.~S.~Petty,
\emph{nfield: A structural analysis engine for fractional
field invariants},
software.
\url{https://github.com/alexspetty/nfield}

\end{thebibliography}

\end{document}
